% ==============================================================================
% Section 5: Non-Parametric Historical Simulation VaR
% Target: Page 7 (Exactly 1 Page)
% ==============================================================================

\section{Non-Parametric Historical Simulation VaR}
\label{sec:historical_simulation_var}

\subsection{Order Statistics and Empirical Percentile Formulation}

Historical Simulation evaluates Value at Risk without imposing a parametric probability distribution on asset returns \citep{beder1995var}. The loss distribution is constructed directly from actual market realizations.

Let $\{R_1, \dots, R_N\}$ denote the sample of $N = 246$ daily historical log returns. Arranging these returns in ascending sequence via order statistics yields:
\begin{equation}
    R_{(1)} \le R_{(2)} \le R_{(3)} \le \dots \le R_{(N-1)} \le R_{(N)}
    \label{eq:order_statistics}
\end{equation}
where $R_{(1)} = -8.47\%$ is the Union Budget excise shock and $R_{(N)} = +5.88\%$ is the sample maximum. For confidence level $\alpha$, the 1-day Historical $\VaR$ is defined by the discrete rank cutoff $k = \lfloor N(1-\alpha) \rfloor$:
\begin{equation}
    \VaR_{\alpha}^{\text{hist}}(1\text{d}) = -V_0 \cdot R_{(\lfloor N(1-\alpha) \rfloor)}
    \label{eq:historical_var_formula}
\end{equation}
Evaluating Equation~\eqref{eq:historical_var_formula} on our $N = 246$ empirical observations ($V_0 = \text{\inr}4,91,625.00$):
\begin{itemize}[leftmargin=1.2em, itemsep=0.5pt, parsep=0pt]
    \item \textbf{1-Day 95\% Conf. ($k=12$):} $R_{(12)} = -2.0945\% \implies \VaR_{0.95}^{\text{hist}}(1\text{d}) = \mathbf{\text{\inr}10,297.08}$ ($2.09\%$).
    \item \textbf{1-Day 99\% Conf. ($k=2$):} $R_{(2)} = -3.9550\% \implies \VaR_{0.99}^{\text{hist}}(1\text{d}) = \mathbf{\text{\inr}19,444.07}$ ($3.96\%$).
    \item \textbf{10-Day Scaled Values ($\times \sqrt{10}$):} $\VaR_{0.95}(10\text{d}) = \text{\inr}32,562.15$ ($6.62\%$), and $\VaR_{0.99}(10\text{d}) = \mathbf{\text{\inr}61,487.89}$ ($12.51\%$).
\end{itemize}

\subsection{The Empirical Fat-Tail Wedge: Parametric vs. Historical}

Comparing Parametric and Historical $\VaR$ reveals a fundamental econometric insight: the \textbf{crossover of risk estimates between 95\% and 99\% confidence levels}.

\begin{table}[!htbp]
    \centering
    \small
    \setlength{\tabcolsep}{5pt}
    \caption{The Fat-Tail Wedge: Parametric vs. Historical Simulation VaR on 1 Lot of ITC Limited ($V_0 = \text{\inr}4,91,625$).}
    \label{tab:fat_tail_wedge}
    \vspace{0.2em}
    \renewcommand{\arraystretch}{1.12}
    \begin{tabularx}{\textwidth}{l r r r >{\raggedright\arraybackslash}X}
        \toprule
        \textbf{Horizon \& Confidence} & \textbf{Parametric} & \textbf{Historical} & \textbf{Divergence (\inr)} & \textbf{Risk Governance Verdict} \\
        \midrule
        1-Day 95\% Confidence & \inr 12,469.05 & \inr 10,297.08 & $-\text{\inr}2,171.97$ & Parametric overstates shoulder risk \\
        1-Day 99\% Confidence & \inr 17,408.01 & \inr 19,444.07 & $+\text{\inr}2,036.06$ & Historical captures empirical fat tail \\
        10-Day 95\% Confidence & \inr 39,430.43 & \inr 32,562.15 & $-\text{\inr}6,868.28$ & Gaussian over-penalizes moderate loss \\
        10-Day 99\% Confidence & \inr 55,048.87 & \inr 61,487.89 & $+\text{\inr}6,439.02$ & \textbf{Parametric understates 10d risk by 10.5\%} \\
        \bottomrule
    \end{tabularx}
\end{table}

At $95\%$ confidence, Parametric $\VaR$ is \textit{higher} than Historical $\VaR$ (\inr 12,469 vs. \inr 10,297) because leptokurtic distributions pull mass from the shoulders into a sharp central peak. However, at $99\%$ confidence, the relationship sharply inverts: \textbf{Historical $\VaR$ exceeds Parametric $\VaR$ by \inr 2,036 on a 1-day basis and by \inr 6,439 on a 10-day basis}. Enforcing a Gaussian bell curve understates the capital required to survive true $99\%$ tail shocks by over $10.5\%$.

\clearpage
